\section{Discussion}

The corrected wireless benchmark does not support a simple claim that one rendering strategy is universally superior. Instead, it shows a trade-off between user-facing load metrics and total client-side energy use. Within this Next.js prototype, SSR reduced median total energy in all three scenarios and did so especially clearly in the \textit{static} and \textit{massive} cases. At the same time, CSR often produced smaller measured transfer sizes and, in the \textit{dynamic} scenario, lower LCP and FCP values. The results are therefore mixed in the sense that different metrics favor different strategies, even though the corrected energy results now lean more strongly toward SSR than the earlier invalidated benchmark suggested.

This pattern is consistent with how the two variants were implemented. In the SSR route, the server resolves the query, loads the scenario-specific inventory, and performs filtering, sorting, pagination, decoration, and summary construction before rendering the page. In the CSR route, the browser receives the page shell first, then fetches raw inventory from the API endpoint after hydration and performs the equivalent transformation on the client. Because both variants share the same application stack, UI, and data model, the comparison isolates where the main computation happens rather than introducing a different framework or toolchain. The observed trade-off is therefore not between unrelated applications, but between server-shifted and client-shifted work within one comparable prototype.

The timing results help explain why CSR can still appear attractive in some scenarios despite its higher energy cost. A smaller initial response and deferred computation can support faster paint-related metrics, particularly when the first visual state is a lightweight shell and the expensive transformation is postponed until after hydration. That interpretation fits the \textit{dynamic} scenario, where CSR achieved lower LCP and FCP than SSR and also lower measured transfer size. However, these earlier paint-related improvements do not necessarily imply lower total client effort. Once the client still has to fetch the inventory and execute the heavy transformation locally, some of the computational cost is shifted to later in the page lifecycle rather than removed altogether.

The corrected wireless run suggests that this deferred client work matters most when the workload grows. In the \textit{massive} scenario, SSR not only reduced median total energy, but also showed a much tighter energy distribution and a lower median JavaScript heap peak than CSR. This is consistent with the idea that moving catalog preparation to the server reduces sustained browser-side processing and memory pressure when the inventory becomes large. The \textit{static} scenario points in the same direction: even with a small dataset, CSR still paid the cost of hydration and a separate inventory request, while the latency advantage that might justify that extra client work was not clearly present in the corrected run. The \textit{dynamic} scenario remained less decisive, indicating that the boundary between ``server work helps'' and ``client work remains acceptable'' is not fixed, but depends on workload and on which metric is prioritized.

Several limitations should temper the interpretation. First, direct energy was estimated from battery current and voltage polling via \texttt{adb dumpsys battery} at 0.25-second intervals rather than from dedicated hardware power-rail counters. This is a practical and defensible method for the available device, but it remains coarser than laboratory-grade power instrumentation. Second, the final conclusions rely on one battery-powered wireless run on a single Samsung Galaxy A53. The randomised order of conditions, fixed screen brightness, and narrow temperature range reduce confounding influences, but they do not eliminate device-specific variability. Third, the energy distributions contained clear high-value runs, particularly in CSR \textit{massive} and CSR \textit{static}, which is why the study emphasizes medians, IQRs, and non-parametric testing. Finally, earlier USB-connected runs were still useful for validating the tooling and the collection pipeline, but USB power could contaminate battery-based energy readings. For that reason, the direct energy interpretation in this paper is based primarily on the corrected battery-powered wireless benchmark rather than on the earlier wired measurements.
